% GENERATED FILE. Edit canonical lessons, then run npm run generate:books.
% canonical-source-hash: fnv1a64:b95e720c8213165d
% canonical-lessons: ZH-C193-tejia, ZH-C193-jiangjia-cheaper, ZH-C193-zhangjia, ZH-C193-jiangjia-haggle, ZH-C193-zhiliang, ZH-R193-review-188-193

\chapter{Prices Up and Down}
\label{ch:zh-prices-up-and-down}
\hlchaptermodality{193}

\begin{chapteropening}
\textbf{By the end of this chapter:} \emph{I can talk about special offers, prices going up and down, haggling and quality.}

Complete the last lesson of chapter 193: I can talk about special offers, prices going up and down, haggling and quality.
\end{chapteropening}

\section[tèjià]{tèjià — special offer; reduced price}

\label{lesson:ZH-C193-tejia}



\begin{quote}
Before the new one: say the Mandarin for to queue, then the Mandarin for to be open to the public.
\end{quote}

\subsection*{You'll want to know: tèjià}
\emph{tèjià} — \textquotedblleft{}special offer; reduced price\textquotedblright{}.

\begin{quote}
\emph{jīntiān shuǐguǒ tèjià} — Fruit is on special offer today.
\end{quote}

Its characters come later; for now it is a word for the ear.

\begin{cousinweb}
\emph{tèjià} is built from \emph{tè} 'special' and \emph{jià} 'price', the \emph{jià} of \emph{jiàqián}, both inherited from Middle Chinese, \emph{dok} and \emph{kaeH} (Guangyun, in Baxter's transcription). The dictionary glosses the first 'a bull' and 'alone'.
\end{cousinweb}

\begin{grammarlens}[title={What you've built}]
One more word for asking for what I want.
\end{grammarlens}

\subsection*{Your turn}
\begin{itemize}
\raggedright
  \item \emph{Say it:} \emph{tèjià}
  \item \emph{Say it:} \emph{tèjià}, once more
  \item \emph{Say it:} \emph{jīntiān shuǐguǒ tèjià} — Fruit is on special offer today.
  \item \emph{Recall:} say the Mandarin for to queue, then the Mandarin for to be open to the public, then say \emph{tèjià} again
\end{itemize}

\begin{tcolorbox}[breakable,colback=teal!4,colframe=teal!35!black,title={Before you move on}]
What does tèjià mean? (\textquotedblleft{}Special offer; reduced price\textquotedblright{}.) Say it once more.
\end{tcolorbox}
\section[jiàngjià]{jiàngjià — to get cheaper; to come down in price}

\label{lesson:ZH-C193-jiangjia-cheaper}



\begin{quote}
Before the new one: say the Mandarin for to be open to the public, then the Mandarin for special offer.
\end{quote}

\subsection*{You'll want to know: jiàngjià}
\emph{jiàngjià} — \textquotedblleft{}to get cheaper; to come down in price\textquotedblright{}.

\begin{quote}
\emph{diànnǎo jiàngjià le} — Computers have got cheaper.
\end{quote}

Its characters come later; for now it is a word for the ear.

\begin{cousinweb}
\emph{jiàngjià} is built from \emph{jiàng} 'to come down' and \emph{jià} 'price', both inherited from Middle Chinese, \emph{kaewngH} and \emph{kaeH} (Guangyun, in Baxter's transcription), which glosses the first 'to go down' and 'to fall'.
\end{cousinweb}

\begin{grammarlens}[title={What you've built}]
One more word for asking for what I want.
\end{grammarlens}

\subsection*{Your turn}
\begin{itemize}
\raggedright
  \item \emph{Say it:} \emph{jiàngjià}
  \item \emph{Say it:} \emph{jiàngjià}, once more
  \item \emph{Say it:} \emph{diànnǎo jiàngjià le} — Computers have got cheaper.
  \item \emph{Recall:} say the Mandarin for to be open to the public, then the Mandarin for special offer, then say \emph{jiàngjià} again
\end{itemize}

\begin{tcolorbox}[breakable,colback=teal!4,colframe=teal!35!black,title={Before you move on}]
What does jiàngjià mean? (\textquotedblleft{}To get cheaper; to come down in price\textquotedblright{}.) Say it once more.
\end{tcolorbox}
\section[zhǎngjià]{zhǎngjià — to get dearer; to go up in price}

\label{lesson:ZH-C193-zhangjia}



\begin{quote}
Before the new one: say the Mandarin for special offer, then the Mandarin for to get cheaper.
\end{quote}

\subsection*{You'll want to know: zhǎngjià}
\emph{zhǎngjià} — \textquotedblleft{}to get dearer; to go up in price\textquotedblright{}.

\begin{quote}
\emph{niúnǎi zhǎngjià le} — Milk has gone up.
\end{quote}

Its characters come later; for now it is a word for the ear.

\begin{cousinweb}
\emph{zhǎngjià} is built from \emph{zhǎng} 'to rise' and \emph{jià} 'price'. The second is inherited from Middle Chinese \emph{kaeH} (Guangyun, in Baxter's transcription). The dictionary glosses the first as water in flood, with a level and a departing tone only, so the rising tone of zhǎng is not the regular outcome, and its history is not settled.
\end{cousinweb}

\begin{grammarlens}[title={What you've built}]
One more word for asking for what I want.
\end{grammarlens}

\subsection*{Your turn}
\begin{itemize}
\raggedright
  \item \emph{Say it:} \emph{zhǎngjià}
  \item \emph{Say it:} \emph{zhǎngjià}, once more
  \item \emph{Say it:} \emph{niúnǎi zhǎngjià le} — Milk has gone up.
  \item \emph{Recall:} say the Mandarin for special offer, then the Mandarin for to get cheaper, then say \emph{zhǎngjià} again
\end{itemize}

\begin{tcolorbox}[breakable,colback=teal!4,colframe=teal!35!black,title={Before you move on}]
What does zhǎngjià mean? (\textquotedblleft{}To get dearer; to go up in price\textquotedblright{}.) Say it once more.
\end{tcolorbox}
\section[jiǎngjià]{jiǎngjià — to haggle}

\label{lesson:ZH-C193-jiangjia-haggle}



\begin{quote}
Before the new one: say the Mandarin for to get cheaper, then the Mandarin for to get dearer.
\end{quote}

\subsection*{You'll want to know: jiǎngjià}
\emph{jiǎngjià} — \textquotedblleft{}to haggle\textquotedblright{}.

\begin{quote}
\emph{zhèr néng jiǎngjià ma} — Can you haggle here?
\end{quote}

Its characters come later; for now it is a word for the ear.

\begin{cousinweb}
\emph{jiǎngjià} is built from \emph{jiǎng} 'to talk over' and \emph{jià} 'price', both inherited from Middle Chinese, \emph{kaewngX} and \emph{kaeH} (Guangyun, in Baxter's transcription), which glosses the first 'to tell' and 'to discuss': to talk the price over.
\end{cousinweb}

\begin{grammarlens}[title={What you've built}]
One more word for asking for what I want.
\end{grammarlens}

\subsection*{Your turn}
\begin{itemize}
\raggedright
  \item \emph{Say it:} \emph{jiǎngjià}
  \item \emph{Say it:} \emph{jiǎngjià}, once more
  \item \emph{Say it:} \emph{zhèr néng jiǎngjià ma} — Can you haggle here?
  \item \emph{Recall:} say the Mandarin for to get cheaper, then the Mandarin for to get dearer, then say \emph{jiǎngjià} again
\end{itemize}

\begin{tcolorbox}[breakable,colback=teal!4,colframe=teal!35!black,title={Before you move on}]
What does jiǎngjià mean? (\textquotedblleft{}To haggle\textquotedblright{}.) Say it once more.
\end{tcolorbox}
\section[zhìliàng]{zhìliàng — quality}

\label{lesson:ZH-C193-zhiliang}



\begin{quote}
Before the new one: say the Mandarin for to get dearer, then the Mandarin for to haggle.
\end{quote}

\subsection*{You'll want to know: zhìliàng}
\emph{zhìliàng} — \textquotedblleft{}quality\textquotedblright{}.

\begin{quote}
\emph{zhège hěn piányi zhìliàng yě hěn hǎo} — It's cheap, and the quality is good too.
\end{quote}

Its characters come later; for now it is a word for the ear.

\begin{cousinweb}
\emph{zhìliàng} is built from \emph{zhì} 'substance' and \emph{liàng} 'a measure', both inherited from Middle Chinese, \emph{tsyit} and \emph{ljangH} (Guangyun, in Baxter's transcription), which glosses the first 'plain' and 'true', and the second as measuring grain.
\end{cousinweb}

\begin{grammarlens}[title={What you've built}]
One more word for asking for what I want.
\end{grammarlens}

\subsection*{Your turn}
\begin{itemize}
\raggedright
  \item \emph{Say it:} \emph{zhìliàng}
  \item \emph{Say it:} \emph{zhìliàng}, once more
  \item \emph{Say it:} \emph{zhège hěn piányi zhìliàng yě hěn hǎo} — It's cheap, and the quality is good too.
  \item \emph{Recall:} say the Mandarin for to get dearer, then the Mandarin for to haggle, then say \emph{zhìliàng} again
\end{itemize}

\begin{tcolorbox}[breakable,colback=teal!4,colframe=teal!35!black,title={Before you move on}]
What does zhìliàng mean? (\textquotedblleft{}Quality\textquotedblright{}.) Say it once more.
\end{tcolorbox}
\section[(dialogue)]{Looking back: from Money to Spend to Prices Up and Down}

\label{lesson:ZH-R193-review-188-193}



\begin{quote}
Say the Mandarin for well-off and to spend money.
\end{quote}

\subsection*{Your turn}
\begin{itemize}
\raggedright
  \item \emph{Say it:} well-off, to spend money, small change, to spend less, enough
  \item \emph{Say it:} day of the month, a form to fill in, to fill in, to sign your name, identity card
  \item \emph{Say it:} bookshop, chemist's, shopping centre, cinema, museum
  \item \emph{Say it:} to come back, to go back there, to go home from abroad, to go abroad, to travel for work
  \item \emph{Say it:} lavatory, lift (in a building), car park, to queue, to be open to the public
  \item \emph{Say it:} special offer, to get cheaper, to get dearer, to haggle, quality
  \item \emph{Read it:} \zh{钱} \zh{号} \zh{店} \zh{回} — each one aloud, then a word that holds it
\end{itemize}

\begin{tcolorbox}[breakable,colback=teal!4,colframe=teal!35!black,title={Before you move on}]
Well-off? (\textbf{\zh{有钱}}, \emph{yǒuqián}.) To spend money? (\emph{huāqián}.) Small change? (\emph{língqián}.) To spend less? (\emph{shěngqián}.) Enough? (\emph{gòu}.) Day of the month? (\textbf{\zh{号}}, \emph{hào}.) A form to fill in? (\emph{biǎogé}.) To fill in? (\emph{tiánxiě}.) To sign your name? (\emph{qiānmíng}.) Identity card? (\emph{shēnfènzhèng}.) Bookshop? (\textbf{\zh{书店}}, \emph{shūdiàn}.) Chemist's? (\emph{yàodiàn}.) Shopping centre? (\emph{shāngchǎng}.) Cinema? (\emph{diànyǐngyuàn}.) Museum? (\emph{bówùguǎn}.) To come back? (\textbf{\zh{回来}}, \emph{huílái}.) To go back there? (\textbf{\zh{回去}}, \emph{huíqù}.) To go home from abroad? (\textbf{\zh{回国}}, \emph{huíguó}.) To go abroad? (\textbf{\zh{出国}}, \emph{chūguó}.) To travel for work? (\emph{chūchāi}.) Lavatory? (\emph{cèsuǒ}.) Lift (in a building)? (\emph{diàntī}.) Car park? (\emph{tíngchēchǎng}.) To queue? (\emph{páiduì}.) To be open to the public? (\emph{kāifàng}.) Special offer? (\emph{tèjià}.) To get cheaper? (\emph{jiàngjià}.) To get dearer? (\emph{zhǎngjià}.) To haggle? (\emph{jiǎngjià}.) Quality? (\emph{zhìliàng}.)
\end{tcolorbox}
